\documentclass[10pt,twocolumn]{article}
\usepackage[T1]{fontenc}
\usepackage[utf8]{inputenc}
\usepackage{lmodern}
\usepackage[margin=1.8cm,top=2.4cm,bottom=2cm,columnsep=0.6cm]{geometry}
\usepackage{fancyhdr}
\usepackage{titlesec}
\usepackage{booktabs}
\usepackage{amsmath,amssymb,amsfonts}
\usepackage{microtype}
\usepackage{xcolor}
\usepackage{mdframed}
\usepackage{parskip}
\usepackage{enumitem}
\usepackage{hyperref}

\definecolor{rulered}{RGB}{180,30,30}
\definecolor{boxgray}{RGB}{245,245,245}
\definecolor{keywordgray}{RGB}{100,100,100}

\pagestyle{fancy}
\fancyhf{}
\fancyhead[L]{\textsc{\small Claude Mag}}
\fancyhead[C]{\small\textit{Machine Learning Research Digest}}
\fancyhead[R]{\small 2026-07-29}
\fancyfoot[C]{\small\thepage}
\renewcommand{\headrulewidth}{0.8pt}

\titleformat{\section}{\normalsize\bfseries\color{rulered}}{}{0pt}{}%
  [\vspace{-4pt}\color{rulered}\rule{\columnwidth}{0.4pt}]
\titlespacing{\section}{0pt}{8pt}{4pt}

\hypersetup{colorlinks=true,urlcolor=rulered,linkcolor=black}

\begin{document}
\setlength{\parindent}{0pt}
\setlength{\parskip}{4pt}

{\small\color{keywordgray}\textbf{Keywords:}
  vision-language-action models \textbullet{}
  robot learning \textbullet{}
  3D pointmaps \textbullet{}
  viewpoint generalization}\par\vspace{4pt}

{\LARGE\bfseries\raggedright
  See like a Robot: Robot-Centric Pointmaps
  for Vision-Language-Action Models}\par\vspace{4pt}

{\large\itshape\raggedright
  Replacing Camera-Frame RGB Images with Robot-Frame 3D Pointmaps
  Closes the Viewpoint Generalization Gap in Large-Scale
  Robot Learning Datasets}\par\vspace{6pt}

{\small\textbf{By} Byungkun Lee, Dongyoon Hwang, Dongjin Kim,
  Hojoon Lee, Minho Park, Jaegul Choo
  (KAIST AI; Holiday Robotics)}\par\vspace{2pt}

{\small\color{keywordgray}arXiv:2607.11498 \textbullet{} July 13, 2026}
\par\vspace{2pt}\noindent{\color{rulered}\rule{\columnwidth}{0.6pt}}\vspace{6pt}

Vision-language-action (VLA) models predict robot actions defined in
the robot's own 3D body frame but observe the world through camera
images defined in the camera's frame---a mismatch that is benign under
fixed cameras but harmful when training datasets aggregate demonstrations
from diverse camera setups. Researchers at KAIST AI and Holiday Robotics
propose \textbf{robot-centric pointmaps}: images whose pixels store the
3D position of scene points in the robot frame, dropping straight into
existing VLA architectures with a single channel-wise concatenation.
Cross-camera evaluation improves by 20 percentage points over RGB-only
baselines.

\begin{mdframed}[backgroundcolor=boxgray,linecolor=rulered,linewidth=1pt,
    innerleftmargin=6pt,innerrightmargin=6pt,
    innertopmargin=6pt,innerbottommargin=6pt]
\textbf{\small AT A GLANCE}\par\vspace{4pt}
\begin{description}[leftmargin=0pt,labelsep=4pt,itemsep=2pt,parsep=0pt,
    font=\small\bfseries]
  \item[Problem:] VLA models observe scenes in the camera frame
    but act in the robot frame; diverse camera setups in large
    datasets make this frame mismatch a primary generalization barrier.
  \item[Why it matters:] Large-scale robot learning requires diverse
    data from many camera rigs; solving viewpoint generalization
    unlocks the full benefit of aggregated datasets.
  \item[Method:] Pointmaps---H$\times$W images with pixel values
    equal to the 3D robot-frame position---are concatenated with RGB
    and passed to an existing VLA with one adapted input layer.
  \item[Result:] Cross-camera success 61.4\% vs.\ 41.2\% for RGB
    only on Open X-Embodiment; $+$20--30pp on real-robot tasks.
  \item[Evidence:] Camera-frame 3D gives only 51.3\% cross-camera
    success, confirming the coordinate system (not just 3D info)
    is the key factor.
\end{description}
\end{mdframed}

\section{The Big Picture}

Scaling robot learning with broad, diverse datasets---analogous to
what large language models did with internet text---is a central goal
in robotics. Datasets like Open X-Embodiment aggregate millions of
demonstrations from dozens of robot platforms and camera setups.
But unlike text, where position in the document is largely irrelevant
to the downstream task, in robot learning the camera viewpoint
fundamentally changes the appearance of the same scene.

A VLA trained on overhead camera demonstrations cannot directly
generalize to wrist-camera deployment because the pixel-to-action
mapping changes completely. If the policy could observe the scene in
the robot's own coordinate frame, viewpoint would become irrelevant:
the spatial relationship between the robot and the object is the same
regardless of which camera captured it.

\section{Why Existing Approaches Fall Short}

\textbf{RGB-only VLAs:} Learn to infer spatial relationships from
monocular cues---perspective foreshortening, shadows, texture gradients.
These cues are camera-dependent and do not generalize across camera
setups without explicit camera-pose conditioning.

\textbf{Camera-frame depth augmentation:} Providing depth in the
camera frame adds metric distance information but not robot-frame
position. The mapping from camera-frame depth to robot-frame coordinates
still requires extrinsic calibration knowledge---which may differ
across demonstrations in a diverse dataset.

\textbf{Camera-pose conditioning:} Explicitly providing camera
extrinsics as input to the policy has been explored, but requires
accurate calibration at deployment and adds architectural complexity.
Policies that memorize pose$\to$action mappings may overfit to
training camera configurations.

\section{Core Mechanism}

A \textbf{robot-centric pointmap} $P \in \mathbb{R}^{H \times W \times 3}$
is an image where each pixel $(u, v)$ stores the 3D coordinate
$(X, Y, Z)$ of the corresponding scene point in the robot's
body frame. The pointmap is constructed at dataset-collection time:

\begin{enumerate}[itemsep=2pt,parsep=0pt]
  \item Backproject each pixel to 3D in the camera frame using depth
    $d_{u,v}$ and intrinsic matrix $K$.
  \item Transform to the robot frame using the known extrinsic
    $T_\text{cam} \to T_\text{robot}$.
  \item Store the result as the pixel value.
\end{enumerate}

The policy receives $[I_\text{RGB}, P] \in \mathbb{R}^{H \times W \times 6}$
as input. The visual encoder's first layer is adapted to accept 6
channels (the 3 new channels initialized to zero weights); all
other weights are kept from pretrained initialization.

\section{Technical Formulation}

For pixel $(u, v)$ with depth $d_{u,v}$, the backprojection is:

\[
  \mathbf{p}_\text{cam} = d_{u,v}\, K^{-1}
    \begin{bmatrix} u \\ v \\ 1 \end{bmatrix}
\]

Transformation to the robot frame via rigid-body transform
$T = (R, \mathbf{t})$:

\[
  \mathbf{p}_\text{robot} = R\,\mathbf{p}_\text{cam} + \mathbf{t}
\]

The pointmap entry is $P_{u,v} = \mathbf{p}_\text{robot}$.
This computation runs at over 60 Hz on standard GPU hardware.

For datasets where depth sensors are unavailable, the authors use
Depth Anything V2 (monocular estimator) followed by scale-and-shift
alignment using known anchor points (e.g., robot end-effector position).
Monocular pointmaps lag behind sensor-depth pointmaps by only 4.2pp
on cross-camera evaluation, showing practical viability without depth
hardware.

\section{Learning or Inference Procedure}

Training follows the standard VLA pipeline---cross-entropy on
action tokens given visual and language input---with the single
change of replacing RGB-only input with 6-channel RGB+pointmap
input. No changes to the loss, optimizer, tokenizer, or action
discretization are required.

At deployment, the robot's camera calibration (intrinsics and
extrinsics relative to its body) is used to compute pointmaps
in real time. The model is never given raw camera-frame depth
or explicit extrinsics as an input token; the coordinate
transformation is baked into the pointmap pixels.

\section{What the Guarantee Says}

If the camera-to-robot transform $T$ is known exactly and the
depth estimate $d_{u,v}$ is accurate, the pointmap provides
exact robot-frame coordinates. Under this assumption, two
demonstrations of the same robot-object configuration from
different cameras produce identical pointmaps---guaranteeing
that the policy input is camera-invariant.

In practice, depth estimation noise and calibration error
introduce per-pixel errors in the pointmap. The empirical
results show that the policy is robust to these errors up to
sensor-depth accuracy levels.

\section{Experimental Findings}

\textbf{Open X-Embodiment cross-camera (task success $\uparrow$):}

\begin{center}\small
\begin{tabular}{lcc}
\toprule
Model & Same-cam. & Cross-cam. \\
\midrule
$\pi_0$ (RGB only) & 74.3 & 41.2 \\
$\pi_0$ + depth (cam.\ frame) & 77.1 & 46.8 \\
$\pi_0$ + pointmap (robot frame) & \textbf{80.6} & \textbf{61.4} \\
\bottomrule
\end{tabular}
\end{center}

\textbf{Real-robot manipulation (50 trials each):}

\begin{center}\small
\begin{tabular}{lcc}
\toprule
Task & RGB only & + Pointmap \\
\midrule
Pick (varied height) & 42\% & 72\% \\
Place (varied surface) & 58\% & 84\% \\
Reach past obstacle & 34\% & 64\% \\
\bottomrule
\end{tabular}
\end{center}

The largest absolute gains appear in tasks requiring precise
distance judgment---picking from varying heights and navigating
reach obstacles---consistent with the hypothesis that robot-frame
metric information is the primary missing signal.

\section{Ablations and Interpretation}

\textbf{Camera-frame vs.\ robot-frame 3D:} Using depth in the camera
frame yields cross-camera success of 46.8\% vs.\ 61.4\% for
robot-frame pointmaps. 3D information helps only when it is
in the robot's native coordinate system.

\textbf{Monocular vs.\ sensor depth:} Sensor depth yields 61.4\%
vs.\ 57.2\% for monocular depth, a 4.2pp gap. Monocular estimation
is a viable substitute in deployment environments where depth sensors
are unavailable.

\textbf{Channel fusion vs.\ early fusion:} Channel concatenation
(6-channel input) outperforms early-layer feature fusion (separate
RGB and 3D branches merged at layer 3) by 2.8pp. Direct concatenation
preserves the ability to use pretrained visual features on RGB while
accessing 3D separately.

\textbf{Effect of dataset diversity:} Training with 10 camera
configurations (vs.\ 2) improves cross-camera accuracy by 8.6pp
for RGB-only but only 2.1pp for pointmap-based models. Pointmaps
substantially reduce the benefit of camera-diverse training data,
confirming that robot-frame grounding eliminates most of the
viewpoint variance.

\section*{Reference}

Byungkun Lee, Dongyoon Hwang, Dongjin Kim, Hojoon Lee, Minho Park,
Jaegul Choo.
\textbf{See like a Robot: Robot-Centric Pointmaps for Vision-Language-Action Models}.
arXiv:2607.11498, July 13, 2026.
\url{https://arxiv.org/abs/2607.11498}

\end{document}
